\section{The intrinsic object and placement}
\label{final:root}

The mathematical object $\Sigma$ is given first in a placed coordinate.
Its placement is recovered exactly, and removing it exposes the
intrinsic decomposition used throughout the paper.

\begin{definition}[The placed presentation of $\Sigma$]
\label{final:C0}
For $\mu>0$, $0<a<1$, put $P=(\mu,a)$, $\gamma=\mu/a$,
$b=-a/\mu$, and $r_*=(1-a)/\mu$. The placed presentation is
\[
 \Sigma_P(r)=\log\!\left(\frac{\mu^2(1+\gamma r)}{\mu+\gamma}\right)-\mu r,
 \qquad r>b.
\]
The intrinsic and placed coordinates, their units, and the parameter
marks are part of these definitions.
\end{definition}

\begin{theorem}[Exact placed reconstruction]
\label{final:C0-placement}
The map $t=\mu r+a$ bijects $(b,\infty)$ with $(0,\infty)$ and gives
\[
 \Sigma_P(r)=\log t-t+\log\frac\mu{1+a}+a.
\]
Inside the placed family, at every point of that domain, placement
is recovered by
\[
 \mu=\sqrt{-\Sigma_P''}-\Sigma_P',\qquad
 \gamma=\frac{\sqrt{-\Sigma_P''}}{1-r\sqrt{-\Sigma_P''}},
 \qquad a=\frac\mu\gamma.
\]
Removing this recovered placement and the value at $r_*$ exposes
the placement-free intrinsic form
\[
 H(t):=\Sigma_P\!\left(\frac{t-a}{\mu}\right)-\Sigma_P(r_*)
      =\log t-t+1,\qquad t>0.
\]
Thus
\begin{equation}
 \Sigma_P(r)=H(\mu r+a)+\log\frac\mu{1+a}+a-1.\label{final:placed}
\end{equation}
Its convex negative $I(t)=t-1-\log t$ and normalized density
$p(t)=te^{-t}$ complete the intrinsic decomposition $(H,I,p)$ of
$\Sigma$, with the exact identities
\begin{align}
 I&=-H,\qquad p=e^{H-1},\qquad H=1+\log p.\label{final:spine}
\end{align}
The closure satisfies $\Sigma_P'(r_*)=0$, $\Sigma_P''(r_*)=-\mu^2$.
The density $p$ has mass one on $(0,\infty)$, whereas
\[
 e^{\Sigma_P(r)}\dd r=\frac{p(t)\dd t}{(1+a)e^{-a}}.
\]
Consequently $e^{\Sigma_P}$ has mass one on $r\ge0$ and mass
$e^a/(1+a)$ on its whole analytic domain.
The presentation $\Sigma_P$ reconstructs both its intrinsic decomposition
$(H,I,p)$ and its placement $P$ exactly; conversely, $(H,I,p)$ together
with $P$ reconstructs $\Sigma_P$ exactly.
\end{theorem}
\begin{proof}
The affine map has positive slope and the stated inverse
$r=(t-a)/\mu$. Substituting $\gamma=\mu/a$ gives
$\mu^2(1+\gamma r)/(\mu+\gamma)=\mu t/(1+a)$, proving
\eqref{final:placed}, with $t=1$ at $r_*$. The logarithm/exponential identities give
\eqref{final:spine}. Integration by parts yields
$\int_a^\infty te^{-t}dt=(1+a)e^{-a}$, including $a=0$.
The change of variables proves both normalization assertions.
Set $q=\gamma/(1+\gamma r)>0$. Differentiation gives
$\Sigma_P'=q-\mu$ and $\Sigma_P''=-q^2$. Solving these equations
gives the recovery formulas, whose denominator is
$(1+\gamma r)^{-1}>0$. At $r_*$ one has $q=\mu$.
\end{proof}

\section{Calibrated analytic reconstruction}
\label{final:core}

\begin{theorem}[Curvature, Riccati and exact flow]
\label{final:C1}
For twice differentiable $F:(0,\infty)\to\R$ with
$F(1)=F'(1)=0$, each of the following conditions identifies $F=H$:
\[
 F''(t)=-t^{-2},\qquad
 F''(t)=-(F'(t)+1)^2.
\]
Equivalently, the marked exact flow
\[
 q(t+s)=\frac{q(t)}{1+s q(t)},\qquad q(1)=1,
\]
for every pair $t,t+s>0$ identifies $q(t)=1/t$; with
$F'=q-1$ and $F(1)=0$ it identifies $F=H$.
More generally $F''(r)=-(r-b)^{-2}$ on $r>b$ identifies
$F(r)=\log(r-b)+Ar+B$, and two independent affine data fix $A,B$.
\end{theorem}
\begin{proof}
Subtract the known logarithmic primitive and use the mean value
theorem twice. For the Riccati equation put $q=F'+1$ and
$W(t)=(tq(t)-1)e^{F(t)+t}$. Direct differentiation using
$q'=-q^2$ gives $W'=0$. The slope anchor gives $W(1)=0$,
so positivity of the exponential gives $tq(t)=1$ everywhere.
Integrating with the value anchor recovers $F$.
For exact flow set $t=1$, $s=v-1$ directly.
All forward assertions follow by differentiating $H$.
Affine perturbations show why the two integration constants are not
fixed by curvature alone.
\end{proof}

\begin{theorem}[Recentring transports its own regularity and anchors]
\label{final:C2}
Let $F$ be differentiable on $(0,\infty)$, with $F''(1)=-1$.
Then
\[
 F(av)-F(a)-aF'(a)(v-1)=F(v)\quad(a,v>0)
\]
holds if and only if $F=H$. No separate affine anchors or global
second-differentiability hypothesis is needed.
\end{theorem}
\begin{proof}
Taking $v=1$ gives $F(1)=0$. Differentiate with respect to $v$ at
one to obtain $F'(1)=0$. At arbitrary $v$ differentiation gives
$a(F'(av)-F'(a))=F'(v)$. The difference quotient at $v=1$ then
gives existence of $F''(a)$ and $a^2F''(a)=F''(1)=-1$.
Theorem~\ref{final:C1} applies. Substitution verifies the converse.
If the scale is not calibrated, every constant multiple of $H$
satisfies the recentering equation.
\end{proof}

\begin{theorem}[Normalized group logarithm and cocycle]
\label{final:C3}
On $D=(-1,\infty)$ define $x\star y=x+y+xy$. This is a
commutative group with identity zero and inverse $-x/(1+x)$.
A function $\ell:D\to\R$ satisfying
\[
 \ell(x\star y)=\ell(x)+\ell(y),\qquad \ell'(0)=1
\]
is uniquely $\ell(x)=\log(1+x)$, where differentiability is
initially required only at zero. Equivalently the differentiable-at-zero
cocycle clause
\[
 h(x\star y)=h(x)+h(y)-xy,\qquad h'(0)=0
\]
identifies $h(x)=\log(1+x)-x=H(1+x)$.
\end{theorem}
\begin{proof}
The identity $1+x\star y=(1+x)(1+y)$ proves closure, associativity,
commutativity and the inverse formula. A homomorphism has $\ell(0)=0$.
For fixed $x$, write $x\star y=x+(1+x)y$, divide its homomorphism
identity by $(1+x)y$, and let $y\to0$. It gives
$\ell'(x)=1/(1+x)$ everywhere. The anchor and the mean value theorem
give the unique primitive. Adding the identity function to the cocycle
turns it into exactly this homomorphism equation. Substitution proves
the forward statements and the cocycle compatibility identity.
\end{proof}

\begin{theorem}[Typed completion equation]
\label{final:C3-completion}
Suppose $h$ is differentiable on $D$, $h'(D)\subset D$ and
$\ell=\mathrm{id}+h$. Then
\[
 \ell(h'(z))=-\ell(z)\quad(z\in D)
\]
holds if and only if, for some $K>0$,
\[
 h(z)=K\log(1+z)-z-\frac K2\log K.
\]
Either anchor $h(0)=0$ or $h'(0)=0$ selects $K=1$.
Injectivity of $\ell$ and twice differentiability of $h$ follow from
the stated domain condition and equation.
\end{theorem}
\begin{proof}
Put $q=h'$. The domain condition gives $\ell'=1+q>0$, so
$\ell$ is strictly increasing and its inverse on its range is
continuous. The equation $q=\ell^{-1}\circ(-\ell)$ makes $q$
continuous. The inverse derivative theorem now makes $q$ continuously
differentiable. Applying the completion equation twice and using
injectivity gives $q(q(z))=z$. Differentiation of the first equation
then gives $(1+z)q'(z)=-(1+q(z))$. Hence
$(1+z)(1+q(z))=K>0$, and integration gives
$h=K\log(1+z)-z+C$. Substitution gives $2C=-K\log K$.
Every displayed function has the requisite derivative range and
satisfies the equation. Its two anchor values are
$-(K/2)\log K$ and $K-1$, each zero exactly when $K=1$.
\end{proof}

\begin{proposition}[Automorphisms, integer series and their calibration]
\label{final:C3-automorphisms}
Every continuous, measurable, monotone or differentiable automorphism
of the fixed group $(D,\star)$ is
\[
 F_c(x)=(1+x)^c-1,\qquad c\in\R\setminus\{0\}.
\]
The calibration $F_c'(0)=1$ selects the identity. Without regularity,
conjugates of arbitrary additive bijections of $\R$ give further
automorphisms. Over a characteristic-zero field the formal
automorphisms have the same formula, with nonzero field coefficient
$c$. The integer series is
$[n]_\star(x)=(1+x)^n-1$; it describes repeated group addition,
and $[m]_\star\circ[n]_\star=[mn]_\star$.
\end{proposition}
\begin{proof}
Conjugate the automorphism by $\log(1+x)$ and $e^u-1$. It becomes
an additive bijection. Each stated regularity condition makes an
additive real function linear: integer and rational linearity follows
algebraically, and continuity follows from measurability or local
boundedness in the other cases. Rational approximation then gives
$u\mapsto cu$, $c\ne0$. A non-linear additive bijection constructed
from a rational basis gives the unrestricted countermodel.
In the formal case, conjugation by the normalized formal logarithm
gives an additive series. Comparing the mixed monomial $u^{n-1}v$
in its additivity identity annihilates every coefficient of degree
$n>1$, since $n$ is invertible. The integer formula follows by the
multiplicative coordinate $1+x$.
\end{proof}

\begin{theorem}[Differentiable Bregman reconstruction]
\label{final:C4}
For a differentiable $\Phi:(0,\infty)\to\R$, let
$B_\Phi(t,s)=\Phi(t)-\Phi(s)-\Phi'(s)(t-s)$.
Simultaneous scale invariance $B_\Phi(ct,cs)=B_\Phi(t,s)$ for all
$c,t,s>0$ holds if and only if
\[
 \Phi(t)=A I(t)+B(t-1)+C.
\]
Thus $\Phi(1)=\Phi'(1)=0$ and
$\Phi'(2)-\Phi'(1)=1/2$ identify $\Phi=I$ without separately
assuming second differentiability or convexity. Alternatively,
for an anchored differentiable $\Phi$ the exact symmetric slice
\[
 B_\Phi(t,1)+B_\Phi(1,t)=t+t^{-1}-2
\]
identifies $\Phi=I$. In particular $B_I(t,s)=I(t/s)$.
\end{theorem}
\begin{proof}
Differentiate scale invariance only in its first argument, taking
$s=1$. For $g=\Phi'-\Phi'(1)$ one obtains
$c(g(ct)-g(c))=g(t)$. Swapping $c,t$ and subtracting gives
$g(c)(1-1/t)=g(t)(1-1/c)$. Fixing $c\ne1$ yields
$g(t)=A(1-1/t)$; integration gives the asserted affine family.
The three calibrations fix its constants. For the symmetric slice,
direct cancellation gives
$(t-1)(\Phi'(t)-\Phi'(1))=t+t^{-1}-2$.
For $t\ne1$ divide by $t-1$ to recover $\Phi'=1-1/t$;
the prescribed derivative at one supplies the remaining point.
Integrate using $\Phi(1)=0$. Direct substitution verifies the forward
identities. Varying $A,B,C$ separately proves the necessity of the
three scalar calibrations within the scale-invariant family.
\end{proof}

\begin{theorem}[Exact Legendre target with either regularity condition]
\label{final:C5}
Extend $I$ by $+\infty$ off $(0,\infty)$. Its convex conjugate is
\[
 I^*(\theta)=\begin{cases}-\log(1-\theta),&\theta<1,\\
 +\infty,&\theta\ge1.\end{cases}
\]
If an extended-real function $F$ has this exact full conjugate and
is either lower semicontinuous on $(0,\infty)$ or convex, then
$F=I$. It is unnecessary to assume both regularity conditions.
Without either, the converse is false.
\end{theorem}
\begin{proof}
For $\theta<1$, the maximizer of $\theta t-I(t)$ is
$t=(1-\theta)^{-1}$, giving the displayed value. The supremum
diverges for $\theta\ge1$. Fenchel's inequality for a candidate
with this conjugate gives $F\ge I^{**}=I$, where the last equality
also follows directly from the just computed maximizers. In particular
$F=+\infty$ off the positive ray. For $t_0>0$, choose
$\theta_0=1-1/t_0$. The function
$(1-\theta_0)t-1-\log t$ is coercive at both endpoints and has
unique minimizer $t_0$. A maximizing sequence in the definition of
$F^*(\theta_0)$ must therefore converge to $t_0$, with its $F$-values
converging to $I(t_0)$. Lower semicontinuity gives the opposite
inequality $F(t_0)\le I(t_0)$. If $F$ is instead convex, these
sequences show its effective domain is dense in the positive ray;
convexity and bracketing by two effective-domain points make it finite
everywhere there. A finite convex function on an open interval is
continuous, so the same argument applies. Finally raise $I(2)$ by
one and leave every other value unchanged. Nearby points approximate
each supremum, so the conjugate is unchanged while the function is
neither convex nor lower semicontinuous at two.
\end{proof}

\begin{theorem}[Self-concordance and the marked Hessian metric]
\label{final:C6}
Let $\Phi\in C^3(0,\infty)$ with $\Phi''>0$,
$\Phi(1)=\Phi'(1)=0$ and $\Phi''(1)=1$. The signed equality
$\Phi'''=-2(\Phi'')^{3/2}$ identifies $\Phi=I$.
The unsigned equality $|\Phi'''|=2(\Phi'')^{3/2}$ also suffices
on the entire positive ray. The corresponding Hessian metric is
$dt^2/t^2$, with distance $|\log(t/s)|$ and
\[
 I(t/s)+I(s/t)=4\sinh^2\!\left(\frac{|\log(t/s)|}{2}\right).
\]
The metric in the marked affine coordinate together with the affine
anchors reconstructs $I$; its bare distance does not fix orientation.
The self-concordance inequality alone does not identify $I$, and
$I$ is not a finite-parameter self-concordant barrier on this domain.
\end{theorem}
\begin{proof}
For the signed equality, $((\Phi'')^{-1/2})'=1$, so calibration gives
$\Phi''=t^{-2}$ and the anchors give $I$. In the unsigned case the
continuous nonzero third derivative has constant sign. The other sign
would give $(\Phi'')^{-1/2}=2-t$, impossible on all $t>0$.
The coordinate $u=\log t$ makes the metric Euclidean, proving the
distance formula; expanding $2\cosh u-2$ proves the symmetrization.
Inversion $t\mapsto1/t$ preserves the distance but reverses the
directed ratio. The quadratic $\Phi(t)=(t-1)^2/2$ satisfies all three
anchors, has positive second derivative, and satisfies the usual
self-concordance inequality since its third derivative is zero,
yet differs from $I$. Moreover,
$(I')^2/I''=(t-1)^2$ is unbounded, refuting a finite barrier parameter.
The affine-equivalent function $-\log t$ does have barrier parameter one.
\end{proof}

\begin{theorem}[Calibrated projective and Schwarzian data]
\label{final:C7}
An injective real-valued function $f$ on $(0,\infty)$ preserving
all cross-ratios and satisfying
$f(1)=0$, $f(2)=-1/2$, $f(3)=-2/3$ equals $1/t-1$.
Consequently $f=F'$ and $F(1)=0$ identify $F=H$.
Alternatively, if $f\in C^3$, $f'\ne0$,
\[
 \frac{f'''}{f'}-\frac32\left(\frac{f''}{f'}\right)^2=0,
 \quad f(1)=0,\quad f'(1)=-1,\quad f''(1)=2,
\]
then the same conclusion holds. Projectivity or zero Schwarzian
without calibration does not select this member.
\end{theorem}
\begin{proof}
With cross-ratio convention
$(x_1-x_3)(x_2-x_4)/((x_1-x_4)(x_2-x_3))$, fix the first three
arguments at $1,2,3$. The cross-ratio is injective fractional-linear
in the fourth argument. Comparing $f$ with $1/t-1$ identifies every
other point; the marks identify the three excluded ones. For the
Schwarzian version put $w=f''/f'$. Its equation is $w'=w^2/2$ with
$w(1)=-2$, whence $w=-2/t$ by ODE uniqueness. Integrate twice
with the two remaining marks to obtain $f'= -t^{-2}$ and
$f=t^{-1}-1$. Integrating the anchored primitive yields $H$.
Without the marks other fractional-linear functions remain.
\end{proof}

\begin{proposition}[Derivative, discrete and local-data boundaries]
\label{final:C8}
The placed derivatives satisfy
\[
 \Sigma_P^{(n)}(r)=(-1)^{n-1}(n-1)!
       \left(\frac\gamma{1+\gamma r}\right)^n,
 \qquad n\ge2.
\]
All orders beyond two add no information once curvature is known.
A lone prescribed derivative of order $n$ identifies a function only
modulo a polynomial of degree at most $n-1$.
For $\rho(r)=C(1+\gamma r)e^{-\mu r}$, $C>0$, its derivative
of order $n\ge1$ has the single zero
$r_n=n/\mu-1/\gamma$. Within this shape family, two consecutive
indexed zeros determine $\mu,\gamma$, and mass determines $C$.
The discrete curvature is
\[
 \Sigma_P(r+1)-2\Sigma_P(r)+\Sigma_P(r-1)
 =\log\left(1-\frac1{(r-b)^2}\right),\quad r-1>b.
\]
It determines lattice values from two starting values, not an arbitrary
real interpolation. A complete germ identifies a global function in a
specified connected real-analytic class, but not in a general smooth class.
\end{proposition}
\begin{proof}
Differentiate the logarithm inductively. Repeated integration proves
the polynomial freedom and Theorem~\ref{final:C1} proves the
second-derivative sufficiency. Leibniz's rule gives
\[
 \rho^{(n)}=C(-\mu)^{n-1}e^{-\mu r}
       \bigl[n\gamma-\mu(1+\gamma r)\bigr],
\]
proving the exact zero. Differences of consecutive zeros recover
$1/\mu$ and then the intercept $-1/\gamma$. Products of the three
logarithmic arguments give the discrete formula; recurrence gives
the lattice-value statement. A nonzero multiple of $\sin(2\pi r)$
preserves all integer samples and changes interpolation. A sufficiently
small compactly supported smooth bump away from an observed point
preserves its entire jet. Two disjoint bumps with cancelling integrals
preserve normalization of a positive density, and can be supported away
from the germ, the mode and endpoints. Analytic continuation instead
uses the identity theorem on the common connected domain.
\end{proof}

\begin{theorem}[Exact all-order derivative-zero sets do not identify $\Sigma$]
\label{final:Z4}
For every real $k>1$ the function
\[
 f_k(t)=(k-1)k^k\frac{t}{(t+k)^{k+1}},\qquad t>0,
\]
is a positive real-analytic probability density, vanishing at both
endpoints, and for every integer $n\ge1$,
\[
 f_k^{(n)}(t)=(-1)^n(k-1)k^k(k)_n
          \frac{t-n}{(t+k)^{k+n+1}}.
\]
Here $(k)_n=k(k+1)\cdots(k+n-1)$. Thus every derivative has
exactly the same single simple positive zero $n$ as $p^{(n)}$,
but $f_k\ne p$. In particular $4t/(t+2)^3$ shows that these
derivative-zero data do not characterize $p$, even among normalized
real-analytic positive densities.
\end{theorem}
\begin{proof}
Substitution $t=ks$ and subtraction of two elementary power integrals
give
\[
 \int_0^\infty \frac{t\,dt}{(t+k)^{k+1}}
 =k^{1-k}\int_0^\infty\frac{s\,ds}{(1+s)^{k+1}}
 =\frac{k^{-k}}{k-1}.
\]
This proves normalization; positivity and real analyticity follow
on the positive ray. The formula is true at order zero if $(k)_0=1$.
Its induction step is the exact identity
\[
 \frac{d}{dt}\frac{t-n}{(t+k)^{k+n+1}}
 =-(k+n)\frac{t-(n+1)}{(t+k)^{k+n+2}}.
\]
Every factor except $t-n$ is nonzero on $t>0$, proving exactness
and simplicity of the zero set, including the derivative sign pattern.
The algebraic tail separates $f_k$ from $te^{-t}$.
The counterexample is not entire and is not globally log-concave:
$(\log f_k)''=-t^{-2}+(k+1)/(t+k)^2>0$ when
$t>1+\sqrt{k+1}$. These additional conditions are not premises of
the refuted characterization. Under $t=\mu r+a$, change variables
and normalize on the selected half-line; multiplication by a positive
constant leaves every placed derivative zero $(n-a)/\mu$ unchanged.
\end{proof}
